\section{Web Agent 实操与挑战}
\subsection{用 Web Agent 规划 GTC 之旅}
在 00:15:41--00:16:39 的演示中，Agent 先问用户“我想去 GTC 2024”，再依次访问 Google、Google Maps，最终找出 San Jose Convention Center——再交由浏览器执行表单与路线导航。该片段展示了 Web Agent 可以在无 API 的场景下依靠页面交互完成任务，也说明它需要不断询问人类、动态 replanning。
\footnote{该片段对应的视频时间：00:15:41--00:16:39。}

\begin{importantbox}{Web Agent 的人机协作边界}
在拿不到 API 的场景下，Web Agent 多靠视觉感知与人类指令循环：它要从浏览器抓取 DOM/文本、再做推理、再执行点击/输入，期间向人类请求补充信息或确认，这就是“方式即能力”的体现。
\end{importantbox}

\subsection{长期规划与观察局限}
Web Agent 的终点往往不在当前页面，因此它必须保持一条内部策略，以免多次点击后才发现走错。除此之外，SRT 00:33:54--00:34:14 指出，即便 Agent 拟定计划，也必须保证执行动作的准确性，否则容易重复环节。

\begin{warningbox}{Planning/Execution 分离的风险}
Web Agent 可能计划得很漂亮，但因为界面变化、按钮文字略有不同、或需要多级验证而失败；此时必须把错误反馈入 tape，让后续的节点学会更鲁棒的执行方式。
\end{warningbox}

\subsection{本章小结}
Web Agent 是“人在 loop”与“Agent 在 loop”共同作用的结果，UI 语义、长期规划与动作执行精度都需要系统化的机制来保障。

